% =============================================================================
% Chapter 6  The Descriptor Reformulation  --  ~9.5 pp  --  THE CORE
% rubric: Project Body (Modelling) -- Model Formulation /30.
% Plan: ThesisPlanning.md §5.6 (headline contribution C).
%
% PURPOSE     Answer question (ii) of sec:pihm-assessment. Written as the methods
%             section of a journal article.
% CRITERION   Formulation band 4, "warranting publication". The body carries four
%             things -- the IDEA, the OBJECT (stated once), the COST (a formal
%             result) and the INTERPRETATION; entry-level formulae, atom-by-atom
%             constructions and proofs go to Appendices B and G.
% SOURCES     Planning.md §3.4 (the state, the block rows, kernel equivalence,
%             bandedness, Neumann on the f'-block, PDEs and coupled fields, the
%             cost of carrying derivatives, block rescaling), §3.7 ("Descriptor:
%             the flat LCU compiled from the term IR"), §6.4; App. A-1, A-2, A-6;
%             claims B-01 to B-07, B-10, B-20, B-21, C-05, C-06;
%             pihm/src/pihm/assemble.py (descriptor mode), ir.py, circuits/lcu.py
%             (built in P5).
% WRITE WHEN  The mathematics now, from Planning.md §3.4, with numbers as \prov;
%             in full at gate G4 (wave W3).
% EQUATION TEST (STYLE_GUIDE.md §2.6): does the argument break if a reader skips
%             it? If not, the equation belongs in an appendix.
% =============================================================================

\subsection{The Idea}
\label{sec:descriptor-idea}
% ~0.75 pp
% LAND, in order:
%   1. One paragraph in words: differentiation of a Chebyshev series is banded
%      if the derivative may live in a neighbouring basis; so carry f', ...,
%      f^(k) as their own blocks and tie them by banded recurrences, rather than
%      substituting the dense G.
%   2. The sparsity figure (F4).
%   3. The spine sentence (ThesisPlanning.md §3.2).
% FIG: fig:sparsity (F4).

The differentiation operator $\G$ is dense. Generally, dense operators are expensive to block-encode and sparse matrices are cheap\todo{Add a reference for this}.

\todo{Remember to define dense vs sparse as part of background}

\placeholderfigure{Dense against banded: the two residuals of the same equation.}{fig:sparsity}{Two sparsity (spy) plots side by side for $\{1,4,4\}$ (R2a) at $n = 4$: left, $\Astd = \G^2 + 4\G + 4I$ with its dense upper triangle; right, $\Adesc$ with its block-banded pattern (recurrence blocks, equation row, padding). Generated from \texttt{pihm.assemble}.}

\subsection{The Banded Factorisation}
\label{sec:descriptor-construction}
% ~0.75 pp
% LAND: G = Bt^{-1} Dt from the recurrence; Bt holds its diagonal and second
%   superdiagonal, Dt its first superdiagonal. RE-CREDIT Olver & Townsend AT THE
%   POINT OF CONSTRUCTION (done below; keep it).
% OFFLOAD: the entries of Bt and Dt -> Appendix B (app:cheb), "the entries of
%   the banded factors, their derivation from the recurrence and the N = 4
%   instance are given in Appendix B".
% Planning.md writes the recurrence as c_r p_r - p_{r+2} = 2(r+1) a_{r+1}; it is
%   eq:cheb-recurrence with the index shifted.

Differentiation of a Chebyshev series is a sparse operation written densely. Let $f = \sum_\ell \fhat{0}_\ell T_\ell$ have derivative $f' = \sum_\ell \fhat{1}_\ell T_\ell$; for $\ell \ge 2$ the coefficient vectors $\fhat{0}$ and $\fhat{1}$ are tied by the recurrence
\begin{equation}
    \fhat{1}_{\ell-1} - \fhat{1}_{\ell+1} = 2\ell\,\fhat{0}_\ell ,
    \label{eq:cheb-recurrence}
\end{equation}
with the $\ell = 1$ row carrying an extra factor of 2 on $\fhat{1}_0$, a normalisation artefact of $T_0$ that \Cref{app:cheb} states in full and derives. Each row relates the coefficient $\fhat{0}_\ell$ to just two coefficients of the derivative. In the normalised Chebyshev basis this is a matrix identity
\begin{equation}
    \Bt\,\fhat{1} = \Dt\,\fhat{0} ,
    \label{eq:banded-factorisation}
\end{equation}
with $\Bt$ banded (its main diagonal and its second superdiagonal), and $\Dt$ carrying a single superdiagonal, such that both hold $\mathcal{O}(N)$ non-zero entries. The dense Chebyshev derivative matrix $\G$ is the operator that sends a function's coefficients directly to those of its derivative, $\fhat{1} = \G\,\fhat{0}$. A simple rearranging of \Cref{eq:banded-factorisation} recovers it as the inverse
\begin{equation}
    \G = \Bt^{-1}\Dt ,
    \label{eq:G-factorisation}
\end{equation}
an identity that holds exactly. The density of $\G$ is due to inverting a banded matrix, and \Cref{eq:banded-factorisation} is the same linear map before that inverse is done. This is the discrete form of the ultraspherical spectral method \cite{olver2013}, in which differentiation becomes sparse once the derivative is allowed to sit in a polynomial basis one order above the function.

\subsection{The Descriptor Residual and Kernel Equivalence}
\label{sec:descriptor-residual}
% ~1.5 pp
% LAND, in order:
%   1. The state w, padded to 2^ceil(log2(k+1)) blocks, and the block rows
%      (kept below; tab:descriptor-rows adds the source and padding rows --
%      merge the two if they read as duplicates).
%   2. thm:kernel-equivalence: exact, given the same M, because Bt is
%      invertible. Proof in the body if it fits in five lines; otherwise
%      Appendix G (app:proof-kernel-equivalence).
%   3. The least-squares ground states differ by row weighting and converge to
%      the same function (restates claim B-02: NOT "bit-identical").
%   4. VOLUNTEER THE COST IN THE SAME BREATH: the Hilbert space grows from N to
%      2^ceil(log2(k+1)) N (the "Cost" paragraph below).
% TAB: tab:descriptor-rows (T5).

Consider the homogeneous order-$k$ linear \gls{ode}
\begin{equation}
  \sum_{j=0}^{k} a_j(x)\, f^{(j)}(x) = 0, \qquad x \in [-1,1],
\end{equation}
where $a_j(x)$ are the (possibly $x$-dependent) coefficients of each $j^\text{th}$ derivative $f^{(j)}$.

We discretise by representing every function appearing in the equation
by its first $N$ coefficients in the normalised Chebyshev basis. Concretely,
let $\fhat{j} \in \mathbb{C}^N$ denote the coefficient vector of $f^{(j)}$, for $j = 0, \dots, k$ and stack these into a single vector,
\begin{equation}
  w = (\fhat{0};\, \fhat{1};\, \dots;\, \fhat{k}).
\end{equation}
Two kinds of linear relations couple the blocks $\fhat{j}$:

\paragraph{Recurrence rows.} The coefficients of consecutive derivatives are related by
\begin{equation}
  -\Dt\, \fhat{j} + \Bt\, \fhat{j+1} = 0, \qquad j = 0, \dots, k-1.
  \label{eq:descriptor-recurrence-rows}
\end{equation}
This gives $k$ rows of the system, each enforcing the derivative relationship between $\fhat{j+1}$ and $\fhat{j}$. These rows are algebraic since they carry no derivative and only fix each $\fhat{j+1}$ as the derivative of $\fhat{j}$, being \Cref{eq:banded-factorisation} applied block-wise and hence equivalent to $\fhat{j+1} = \G \fhat{j}$ on $\ker \Adesc$.

\paragraph{Differential row.} The differential equation itself becomes one
further row, obtained by multiplying each coefficient vector $\fhat{j}$ by
the (possibly $x$-dependent) coefficient $a_j(x)$ and summing:
\begin{equation}
  \sum_{j=0}^{k} \M_{a_j}\, \fhat{j} = 0,
  \label{eq:descriptor-ode-row}
\end{equation}
where $\M_{a_j}$ is the matrix representing multiplication by $a_j(x)$ in
coefficient space (reducing to $a_j I_N$ when $a_j$ is constant). This single row is the equation itself: we carry the intermediate derivatives as their own unknown coefficient vectors within $w$, rather than representing them via $\fhat{j} = \G^j\,\fhat{0}$, so that \Cref{eq:G-factorisation} (the only dense step) never appears.

Together these $k+1$ block rows ($k$ recurrence rows and one
differential row) define a linear system for $w$ and transform the original \gls{ode} problem into a \gls{dae} system.

\begin{definition}[Descriptor form residual operator]
\label{def:descriptor-residual-operator}
Collecting all block rows, the descriptor form chebyshev residual operator
$\Adesc$ is the $(k+1)N \times (k+1)N$ block matrix
\begin{equation}
\Adesc =
\begin{pmatrix}
-\Dt & \Bt &        &        &          \\
          & -\Dt & \Bt &      &          \\
          &          & \ddots & \ddots &          \\
          &          &        & -\Dt & \Bt \\
\M_{a_0}   & \M_{a_1}  & \cdots & \cdots & \M_{a_k}
\end{pmatrix},
\end{equation}
where the first $k$ rows are the recurrence rows and the $(k+1)$-th row is
the differential row, so that the whole system is compactly written as
\begin{equation}
  \Adesc\, w = 0.
\end{equation}
\end{definition}

\paragraph{A first-order example.}
For $f' + 2f = 0$ the state is $w = (\fhat{0};\, \fhat{1})$, and the two facts relating its blocks are the recurrence $\Bt\fhat{1} - \Dt\fhat{0} = 0$ and the equation $2\fhat{0} + \fhat{1} = 0$, giving
\begin{equation}
    \Adesc =
    \begin{pmatrix}
        -\Dt & \Bt \\[2pt]
        2 I_N & I_N
    \end{pmatrix},
    \qquad w = \begin{pmatrix} \fhat{0} \\ \fhat{1} \end{pmatrix} .
    \label{eq:descriptor-first-order}
\end{equation}
Every block is $\Bt$, $\Dt$ or a multiple of the identity, so $\Adesc$ is banded, against the dense $\G + 2I$ that substituting $\fhat{1} = \G \fhat{0}$ would produce. The fully populated $N = 4$ instance of \Cref{eq:descriptor-first-order}, the entries of $\Bt$ and $\Dt$, and the exactness of \Cref{eq:G-factorisation}, are given in \Cref{app:cheb}.

\begin{table}[htbp]
    \centering
    \caption{The block rows of the descriptor residual $\Adesc$ for an order-$k$ equation.}
    \label{tab:descriptor-rows}
    \small
    \begin{tabular}{@{}l l@{}}
        \toprule
        Block row & Equation \\
        \midrule
        $j = 0, \dots, k-1$ (recurrence) & $-\Dt\,\fhat{j} + \Bt\,\fhat{j+1} = 0$ \\
        $k$ (the equation) & $\sum_{j} \M_{a_j}\,\fhat{j} - \M_r\,\Dzero\,\fhat{0} = 0$ \\
        $b > k$ (padding) & $c_{\mathrm{pad}}\,\fhat{b} = 0$ \\
        \bottomrule
    \end{tabular}
\end{table}

\begin{theorem}[Kernel equivalence]
\label{thm:kernel-equivalence}
% STATE: if both forms use the same multiplication operators M, the descriptor
%   kernel is exactly the image of the standard-form kernel under the map
%   p -> (p, Gp, ..., G^k p). The proof needs only that Bt is upper triangular
%   with a non-zero diagonal, hence invertible.
\begin{equation}
    \ker \Adesc = \bigl\{\, (p,\ \G p,\ \dots,\ \G^{k} p) \;:\; p \in \ker \Astd \,\bigr\} .
    \label{eq:kernel-equivalence}
\end{equation}
\end{theorem}

\paragraph{Cost.}
The reformulation trades a higher dimensional Hilbert space (and with it a higher qubit count) for a sparser residual operator. $w$ has $(k+1)N$ components at the classical, unpadded level; realising the field register on $\lceil\log_2(k+1)\rceil$ qubits pads this to $2^{\lceil\log_2(k+1)\rceil}N$ (\Cref{sec:app-variable-coeff}), which is what should be compared against the $N$ of the standard form. Each block of $\Adesc$ is $\Bt$, $\Dt$, or a banded multiplication operator (bandwidth $q$ for a degree-$q$ polynomial coefficient; \Cref{sec:app-variable-coeff}), so $\Adesc$ carries $\mathcal{O}(1)$ occupied diagonals per block, whereas the standard-form residual $\Astd = \sum_j a_j \G^j$ is dense with $\Theta(N^2)$ non-zero entries. The $\fhat{0}$-block of the recovered null vector is the Chebyshev coefficient vector the standard form solves for, agreeing with it to within numerical precision, quantified in \Cref{sec:results}. Carrying the intermediate derivatives instead of eliminating them converts the dense physics-informed residual into a banded one, and every later stage of the pipeline inherits that structure.
% REVISE: (1) "agreeing to within numerical precision" -> state it through
%   thm:kernel-equivalence: exact kernels coincide; least-squares ground states
%   differ by row weighting and converge to the same function (claim B-02).
%   (2) Update the closing sentence to the spine of ThesisPlanning.md §3.2
%   (Theta(N^8) -> Theta(N^2), kernel unchanged).

\subsection{Constraints, Partial Differential Equations and Coupled Fields}
\label{sec:descriptor-constraints}
% ~1 pp
% LAND, in order:
%   1. The shared constraint layer (tab:constraint-kinds) acts on chosen blocks.
%      Neumann sits on the f'-block: the same atom as Dirichlet, with no G --
%      this retires claim B-10 ("Neumann costs Theta(N^2) more").
%   2. PDEs: a shared f-block plus one derivative chain per axis. Coupled
%      fields: one chain per field.
%   3. Datum slices for initial and boundary FUNCTIONS. These are available to
%      the standard form too, so "the descriptor avoids doubling for linear
%      PDEs" is NOT claimed (B-20).

\subsection{Block Rescaling}
\label{sec:descriptor-rescaling}
% ~0.75 pp
% LAND, in order:
%   1. Because ||G|| = Theta(N^2), derivative blocks dominate the state's norm
%      and the operator's scaling: Delta = 7.5e-6 for {1,5,400} unrescaled
%      (\prov; Planning.md A-1).
%   2. The remedy (eq:block-rescaling) is a diagonal on the field register, at
%      Theta(n_f) gates.
%   3. Why general equilibration is excluded from circuit paths: Theta(P) gates
%      (validity issue V-08).
%   4. zeta is problem-derived (zeta = omega for R2c), not fitted (Planning.md
%      §14 item 4).
% Notation: Planning.md writes this constant lambda; the thesis uses zeta,
%   because lambda is taken by eigenvalues (Delta = lambda_1 - lambda_0).

\begin{equation}
    \fhat{m} \;\longmapsto\; \zeta^{-m}\,\fhat{m},
    \qquad m = 0, \dots, k .
    \label{eq:block-rescaling}
\end{equation}

\subsection{The Block Encoding: One Flat Linear Combination}
\label{sec:descriptor-block-encoding}
% ~1.5 pp
% LAND, in order:
%   1. Every block is a short sum of SHIFT^s x diag terms; the diagonal of Dt is
%      linear in the index (eq:diag-index); nearly constant diagonals use
%      reflections R_m = I - 2|m><m|. So A_desc is ONE flat LCU of O(n) terms
%      (eq:descriptor-lcu).
%   2. The term IR compiles to PREPARE/SELECT/PREPARE^dagger: the quantikz figure
%      (F5) and the algorithm listing (alg:ir-to-lcu), followed by a
%      line-referenced walkthrough (Snow's model: "Inputs and outputs",
%      "PREPARE (line 2)", "SELECT (lines 3-5)").
%   3. Rank-one rows via the feature map; NO dense state preparation anywhere
%      (implementation issue I-01, fixed).
% FIG: fig:descriptor-lcu (F5). TAB (optional): tab:descriptor-atoms.
% The listing below is a placeholder: replace its lines with the compiler's
%   real steps (pihm/src/pihm/ir.py, circuits/lcu.py) once P5 has built them.

\begin{equation}
    \operatorname{diag}(k)_{k=0}^{N-1} = \sum_{j=0}^{n-1} 2^{j}\,\frac{I - Z_j}{2},
    \qquad
    R_m = I - 2\ket{m}\bra{m} .
    \label{eq:diag-index}
\end{equation}
\begin{equation}
    \Adesc = \sum_{t=1}^{L} c_t\,U_t,
    \qquad
    U_t = (\text{field word}) \otimes \mathrm{SHIFT}^{s_t}\,P_t,
    \qquad
    \alpha_A = \sum_{t=1}^{L} |c_t| .
    \label{eq:descriptor-lcu}
\end{equation}

\placeholderfigure{The descriptor block encoding: one PREPARE/SELECT/PREPARE$^\dagger$ over the whole residual.}{fig:descriptor-lcu}{Quantikz circuit: selector register (PREPARE loads $\sqrt{|c_t|/\alpha_A}$), field register ($n_f$ qubits), system register ($n$ qubits); SELECT applies each term's field word, cyclic shift $\mathrm{SHIFT}^{s_t}$ (multi-controlled-$X$ ladder) and $Z$-string or reflection $P_t$; PREPARE$^\dagger$; post-selection on the selector.}

\begin{algorithm}[htbp]
    \caption{Compiling the term IR of $\Adesc$ into PREPARE and SELECT.}
    \label{alg:ir-to-lcu}
    \KwIn{the term IR $\{(c_t, U_t)\}_{t=1}^{L}$ of $\Adesc$ (\Cref{eq:descriptor-lcu})}
    \KwOut{PREPARE and SELECT circuits; the subnormalisation $\alpha_A$}
    $\alpha_A \leftarrow \sum_{t} |c_t|$\;
    PREPARE: $\ket{0} \mapsto \sum_{t} \sqrt{|c_t|/\alpha_A}\,\ket{t}$\;
    \ForEach{term $t$}{
        SELECT: controlled on $\ket{t}$, apply $\operatorname{sgn}(c_t)\,U_t$\;
    }
    \Return{PREPARE, SELECT, $\alpha_A$}\;
\end{algorithm}

\begin{table}[htbp]
    \centering
    \caption{The atoms of the descriptor block encoding.}
    \label{tab:descriptor-atoms}
    \small
    \begin{tabular}{@{}l c c c c@{}}
        \toprule
        Atom & Arity & Gates & Ancilla & Verified residual \\
        \midrule
        Cyclic shift $\mathrm{SHIFT}^{s}$ & & & & \\
        $Z$-string diagonal & & & & \\
        Reflection $R_m$ & & & & \\
        Feature-map rank-one row & & & & \\
        Block-rescaling diagonal & & & & \\
        \bottomrule
    \end{tabular}
\end{table}

\subsection{Cost of the Descriptor Encoding}
\label{sec:descriptor-cost}
% ~1.25 pp
% LAND, in order:
%   1. prop:descriptor-cost (claims B-04, C-06): alpha_A = 2N + O(1), with
%      alpha_A / ||A_desc|| -> 1 (measured 2.29 -> 1.002 over n = 2-12, \prov);
%      ancilla ceil(log2(n+4)) + 6 for R2's family; O(n^2) gates.
%   2. With Omega(log L) for exact LCU (sec:bg-block-encoding), the ancilla count
%      is order-optimal AMONG EXACT LCU ENCODINGS. Hedge exactly that far.
%   3. Two paragraphs of OPERATIONAL INTERPRETATION in the body: what 10
%      ancilla at n = 12 means next to the standard form's Omega(n); what
%      alpha -> ||A|| means for the filter.
% OFFLOAD: proof -> Appendix G (app:proof-descriptor-cost).

\begin{proposition}[Cost of the descriptor encoding]
\label{prop:descriptor-cost}
% STATE: for the second-order family R2 (and analogously for any fixed order),
%   the flat LCU of eq:descriptor-lcu is an exact block encoding of A_desc with
%   the subnormalisation, ancilla count and gate count below.
\begin{equation}
    \alpha_A = 2N + O(1),
    \qquad
    \frac{\alpha_A}{\|\Adesc\|_2} \to 1,
    \qquad
    \lceil \log_2(n+4) \rceil + 6\ \text{ancilla},
    \qquad
    O(n^2)\ \text{gates} .
    \label{eq:descriptor-cost}
\end{equation}
\end{proposition}

\subsection{Consequences for the Spectrum and the Filter}
\label{sec:descriptor-spectrum}
% ~1.25 pp
% LAND, in order:
%   1. ||H_desc|| = Theta(N^2); alpha_H / ||H_desc|| -> 1 (5.3 -> 1.05, \prov).
%   2. Delta is problem-dependent: flat at 0.454 for {1,4,4}; 7.5e-6 for
%      {1,5,400} before rescaling (\prov).
%   3. Hence, by eq:degree, d = Theta~(N) against Theta~(N^4) (eq:standard-wall)
%      -- eq:descriptor-degree (claim C-05).
%   4. The relative gap, 7.0e-6 against 3.0e-15 at n = 7 for {1,4,4} (\prov),
%      stated PRECISELY: the standard form's ground state is below double-
%      precision resolution by eigh, NOT by SVD. The consequence is the degree,
%      not unresolvability (restates claims B-07 and B-23).
%   5. Point to sec:results-scaling for the measured curves.

\begin{equation}
    \|\Hdesc\| = \Theta(N^2),
    \qquad
    \frac{\alpha_H}{\|\Hdesc\|} \to 1
    \quad\Longrightarrow\quad
    d = \tilde\Theta(N) .
    \label{eq:descriptor-degree}
\end{equation}

\subsection{Summary}
\label{sec:descriptor-summary}
% ~0.25 pp
% LAND: four numbered takeaways, then the spine sentence (ThesisPlanning.md §3.2).

% CHECKLIST (marker)
%   1. Could a referee re-derive A_desc from the body plus Appendices B and G?
%   2. Is the Hilbert-space cost volunteered alongside the gains?
%   3. Is the ultraspherical antecedent credited twice (sec:bg-ultraspherical
%      and sec:descriptor-construction)?
%   4. Is the cost result a proposition AND interpreted?
%   5. Is the ancilla claim O(log n), with optimality hedged to exact LCU?
%   6. Does the gap discussion distinguish eigh failure from true
%      unresolvability?
%   7. Does every body equation pass the equation test?
% TRAPS  "O(1) ancilla"; "bit-identical a-block"; "only the descriptor can
%        prepare nonlinear solutions"; equilibration numbers quoted without their
%        Theta(P) gate cost.
